\documentclass[11pt,oneside]{book}
\usepackage[T1]{fontenc}
\usepackage{lmodern}
\usepackage{principia}
\usepackage{microtype}
\usepackage[hidelinks]{hyperref}

\title{A Working Notebook in Peanese Notation}
\author{A. Researcher}
\date{\today}

\begin{document}
\frontmatter
\maketitle
\tableofcontents

\chapter{How to use this starter}
This starter is a compact book-shaped example for the \texttt{principia}
package. Keep ordinary exposition in the surrounding chapters and use the
package commands when a proposition needs the distinctive notation of
Whitehead and Russell's \emph{Principia Mathematica}.

\mainmatter
\chapter{A small formal system}
\section{Propositions and scope}
The Peanese connectives are ordinary LaTeX math material, so they can be
introduced inline: $p \pmimp q$ and $p \pmiff q$. Square scope dots are
available in several lengths, for example
\[
  \pmnot p \pmor q,\qquad
  p \pmandd q \pmimp r,\qquad
  p \pmanddd q \pmiff r.
\]
The starred proposition marker and definition notation make a useful compact
house style for numbered arguments:
\[
  \pmsn{2}{1}\; p \pmimp p \pmdf
\]
\section{A reusable argument pattern}
Let $P$ and $Q$ be propositional functions. A substitution can be set out
without drawing a custom symbol each time:
\[
  \pmsub{P x}{P a},\qquad
  \pmsubbb{P x}{P a}{Q x}{Q a}{R x}{R a}.
\]
Use the macros in \texttt{principia.sty} as the source of truth for the
notation; this document intentionally demonstrates a small, readable subset.

\chapter{Notes for an author}
\begin{itemize}
  \item Keep logical expressions in math mode so spacing remains stable.
  \item Put long deductions in displayed equations and explain each step in
        prose immediately before or after it.
  \item Extend the starter by adding chapters and by selecting the volume-
        specific commands documented in the upstream package manual.
\end{itemize}

\backmatter
\chapter*{Colophon}
This starter uses the LPPL-1.3c \texttt{principia} package by Landon D. C.
Elkind. No upstream example prose, images, or logo are bundled.
\end{document}
